\documentclass[../../../include/open-logic-section]{subfiles}

\begin{document}

\olfileid{his}{set}{pathology}
\olsection{غيرمعمولي جوړښتونه}

خو د \emph{حد} تعريف د ځينو داسې جوړښتونو اجازه ورکړه چې تر ډېره «غيرمعمولي» بلل کېدل۔

د 1830مو کلونو په شاوخوا کښې بولڅانو داسې يوه تابع وموندله چې \emph{په هره نقطه کښې پيوسته} وه، خو \emph{په هېڅ نقطه کښې يې مشتق نۀ درلود}۔ (له بده مرغه، بولڅانو دا کار هېڅکله خپور نۀ کړ؛ رياضيپوهان په لومړي ځل په 1872 کښې له دې نظر سره اشنا شول، کله چې وايرشټراس په خپلواکه توګه همدغه نظر وموند۔)\footnote{دا تاريخ د ويکيپېډيا د \href{http://en.wikipedia.org/wiki/Weierstrass_function}{وايرشټراس د تابع} په مقاله کښې په ډېر مفصلو لمنليکونو کښې بيان شوے دے۔} لږ تر لږه دا خبره ډېره حيرانوونکې وه۔ داسې تابعې موندل اسانه دي، لکه $|x|$، چې په هره نقطه کښې پيوستې وي خو په يوه ځانګړې نقطه کښې مشتق نۀ لري۔ خو داسې تابع چې په هره نقطه کښې پيوسته وي او په \emph{هېڅ} نقطه کښې مشتق نۀ لري، ډېر بل ډول څيز دے۔ يوه شېبه فکر وکړئ چې د داسې تابع د رسمولو هڅه به څنګه کوئ۔ د پيوستون لپاره بايد دا د پاڼې څخه د قلم پورته کولو پرته رسم کړے شئ؛ خو د دې لپاره چې په هېڅ نقطه کښې مشتق ونۀ لري، بايد د خپل قلم لورے په پرله‌پسې توګه ناڅاپه بدل کړئ۔

نور «غيرمعمولي جوړښتونه» هم راپسې راغلل۔ د 1874 د جنورۍ په 5مه نېټه کانتور ډېډېکنډ ته يو ليک وليکه او دا مسئله يې وړاندې کړه:
\begin{quote}
آيا يوه سطحه (مثلاً يوه مربع چې څنډه يې هم پکې شامله وي) له يوه خط سره (مثلاً يو مستقيم خط چې دواړه پای‌ټکي يې هم پکې شامل وي) يو پر يو داسې برابرېدای شي چې د سطحې هرې نقطې ته د خط يوه نقطه برابره وي، او برعکس، د خط هرې نقطې ته د سطحې يوه نقطه برابره وي؟

اوس هم راته داسې ښکاري چې د دې پوښتنې ځواب ډېر ګران دے؛ سره له دې چې دلته هم انسان دومره اړ کېږي چې \emph{نه} ووايي، چې غواړي ثبوت يې تقريباً بې‌اړتيا وګڼي۔ [په \citealt{Gouvea2011} کښې نقل شوے دے۔]
\end{quote}
خو په 1877 کښې کانتور ثابته کړه چې د هغه پخوانے ګومان ناسم و۔ په رښتيا، يو خط او يوه مربع عين يو هومره نقطې لري۔ د 1877 د جون په 29مه نېټه يې ډېډېکنډ ته وليکل: «\emph{je le vois, mais je ne le crois pas}»؛ يعنې «وينم يې، خو باور پرې نۀ کوم»۔ د رياضياتو په «رواجي تاريخ» کښې دا خبره ډېر ځله د دې نښې په توګه اخيستل کېږي چې دا نوې نتيجې د هغه وخت د رياضيپوهانو لپاره څومره \emph{په لفظي معنا د باور وړ نۀ وې}۔ (دا ليکونه په \citet{Gouvea2011} کښې وړاندې شوي، او په \olref[his][set][mythology]{sec} کښې ورته بېرته راځو۔ د کانتور د ثبوت خاکه په \olref[his][set][cantorplane]{sec} کښې شته۔)

د کانتور له نتيجې څخه په الهام، پېانو په دې غور پيل کړ چې آيا يو خط پر يوه مستوي باندې \emph{په نرمۍ} تصوير کېدای شي۔ دا به يو \emph{ځای ډکوونکے منحني خط} وي۔ په \citeyear{Peano1890} کښې پېانو همداسې يو منحني خط جوړ کړ۔ \emph{سرچينه‌يي يادونه OLSTH-049: دلته «په نرمۍ» د پيوسته تصوير غيررسمي بيان دے؛ دا د پيوسته مشتق لرلو تخصصي شرط نۀ دے۔ له يوې تړلې متناهي عددي وقفې څخه د مربع پر ټولو نقطو پيوسته تصوير شته، خو د پيوسته مشتق لرونکے داسې تصوير مربع نۀ شي ډکولے۔ اصلي عبارت ساتل شوے او د دواړو معناوو توپير څرګند شو۔} دا خبره په رښتيا د هندسي وجدان خلاف ښکاري: اقليدس خط «بې پلنوالي اوږدوالے» تعريف کړے و (لومړے کتاب، تعريف 2)، خو پېانو ښودلې وه چې يو خط په مناسبه توګه تاو راتاو کولو سره د هغه اوږدوالے په پلنوالي بدلېدای شي۔ په \citeyear{Hilbert1891} کښې هېلبرټ يو څه زيات له وجدان سره سم ځای ډکوونکے منحني خط بيان کړ او د هغه د څرګندولو لپاره يې ځينې انځورونه هم وړاندې کړل۔ دا منحني خط په پرله‌پسې پړاوونو کښې جوړېږي، او دلته د جوړښت لومړني شپږ پړاوونه دي:
\begin{center}
\begin{tikzpicture}
\begin{scope}[xshift=20pt, yshift=20pt]
	\draw [oldiagcolorC, l-system={Hilbert curve, step=40pt, angle=90, axiom=L, order=1}] lindenmayer system;
\end{scope}
\begin{scope}[xshift=110pt, yshift=10pt]
	\draw [oldiagcolorC, l-system={Hilbert curve, step=20pt, angle=90, axiom=L, order=2}] lindenmayer system;
\end{scope}
\begin{scope}[xshift=205pt, yshift=5pt]
	\draw [oldiagcolorC, l-system={Hilbert curve, step=10pt, angle=90, axiom=L, order=3}] lindenmayer system;
\end{scope}
\begin{scope}[xshift=2.5pt, yshift=-97.5pt]
	\draw [oldiagcolorC, l-system={Hilbert curve, step=5pt, angle=90, axiom=L, order=4}] lindenmayer system;
\end{scope}
\begin{scope}[xshift=101.25pt, yshift=-98.75pt]
	\draw [oldiagcolorC, l-system={Hilbert curve, step=2.5pt, angle=90, axiom=L, order=5}] lindenmayer system;
\end{scope}
\begin{scope}[xshift=200.625pt, yshift=-99.375pt]
	\draw [oldiagcolorC, l-system={Hilbert curve, step=1.25pt, angle=90, axiom=L, order=6}] lindenmayer system;
\end{scope}
\draw[gray] (0pt,0pt) rectangle (80pt, 80pt);
\draw[gray] (100pt,0pt) rectangle (180pt, 80pt);
\draw[gray] (200pt,0pt) rectangle (280pt, 80pt);
\draw[gray] (0pt,-100pt) rectangle (80pt, -20pt);
\draw[gray] (100pt,-100pt) rectangle (180pt, -20pt);
\draw[gray] (200pt,-100pt) rectangle (280pt, -20pt);
\end{tikzpicture}  
\end{center}
په حد کښې — هغه مفهوم چې تر دې وخته يې کره تعريف ترلاسه کړے و — ټوله مربع په بشپړ ډول په !!{colorC} رنګ ډکېږي۔ د يادونې وړ دا هم ده چې د هېلبرټ منحني خط په هره نقطه کښې پيوسته دے، خو په هېڅ نقطه کښې مشتق نۀ لري؛ د وجدان له مخې يې وجه دا ده چې په نامتناهي حد کښې تابع په هره شېبه کښې ناڅاپه خپل لورے بدلوي۔ (د هېلبرټ د جوړښت لا مفصله خاکه به په \olref[his][set][hilbertcurve]{sec} کښې وړاندې کړو۔)

د ښه يا بد لپاره، دا «غيرمعمولي» هندسي جوړښتونه پر هندسي وجدان د تکيه کولو د شک لپاره د دليل په توګه واخيستل شول۔ دا د رياضياتو د يوې نوې طريقې لپاره د \emph{تبليغاتو} تر کچې ورنژدې شول؛ داسې طريقه چې په پای کښې به سټ نظريې ته ورسېږي۔ وروسته، د دې مضمون په اسطوري بيانونو کښې، دا خبره بيا بيا وشوه چې دا نتيجې هم پوره کره وې او هم پوره ټکان ورکوونکې۔ ځکه نو دوه موخې يې ترسره کولې: پر هندسي وجدان د تکيه کولو پر ضد خبردارے، او د فکر د نويو طريقو د حاصل ورکولو څرګندونه۔

\end{document}
